% 陈冠宇 — Curriculum Vitae
% Compile: xelatex cv.tex (run twice for refs)

\documentclass[11pt,a4paper]{article}

% --- Encoding & fonts ---
\usepackage{fontspec}
\usepackage{xeCJK}

% English (Latin) font: classic serif close to the web CV
\setmainfont{EB Garamond}[
  UprightFont = *,
  ItalicFont = *,
  BoldFont = *,
  Ligatures = TeX
]

% CJK fonts: use Fandol (bundled with TeX Live, static, no VF gotchas).
% Song = 宋体 (serif body), Hei = 黑体 (sans for section heads), Kai = 楷体 (italic-ish).
\setCJKmainfont{FandolSong}
\setCJKsansfont{FandolHei}
\setCJKfamilyfont{kai}{FandolKai}
\xeCJKsetup{CJKecglue={\hskip 0.15em plus 0.05em}}

% --- Layout ---
\usepackage[a4paper,top=0.85in,bottom=0.85in,left=0.95in,right=0.95in]{geometry}
\usepackage{xcolor}
\usepackage{enumitem}
\usepackage[hidelinks,unicode]{hyperref}
\usepackage{calc}

% --- Colors (match the site palette) ---
% Use named SVG colors so xcolor doesn't need the HTML model.
\definecolor{rulecol}{rgb}{0.792,0.835,0.882}
\definecolor{subdued}{rgb}{0.392,0.455,0.545}
\definecolor{bodycol}{rgb}{0.118,0.161,0.231}
\definecolor{accentcol}{rgb}{0.059,0.090,0.165}
\color{bodycol}

% --- Section heading style: small caps + rule, like the web CV ---
\usepackage{titlesec}
\titleformat{\section}
  {\sffamily\fontsize{10.5pt}{12pt}\selectfont\bfseries\color{accentcol}\MakeUppercase}
  {\makebox[1.6em][l]{\rmfamily\fontsize{9pt}{11pt}\selectfont\color{subdued}\thesection.}}{0pt}
  {}
  [\vspace{-2pt}{\color{rulecol}\rule{\linewidth}{0.4pt}}%
   \rmfamily\mdseries\upshape\color{bodycol}]
\titlespacing*{\section}{0pt}{1.3em}{0.55em}

% --- Custom entry helper ---
% Args: 1=date, 2=title, 3=subtitle, 4=bullet items (or empty)
\newcommand{\cventry}[4]{%
  \par\noindent
  \makebox[\linewidth][s]{%
    \sffamily\bfseries\color{accentcol}#2%
    \hfill
    \rmfamily\itshape\color{subdued}\fontsize{10pt}{12pt}\selectfont #1%
  }\par
  \vspace{0.15em}
  {\itshape\color{subdued}#3}\par
  \ifx&#4&\else
    \vspace{0.2em}
    \begin{itemize}[leftmargin=1.25em,labelsep=0.4em,itemsep=0.05em,topsep=0pt,parsep=0pt]
      #4
    \end{itemize}
  \fi
  \vspace{0.55em}
}

% --- Skill row helper ---
% Args: 1=key (uppercased, gray, sans), 2=value (body, wraps if long)
\newcommand{\cvskill}[2]{%
  \par\noindent
  \makebox[8em][l]{\sffamily\fontsize{9.5pt}{11pt}\selectfont\color{subdued}\MakeUppercase{#1}}%
  \color{bodycol}#2%
  \par
  \vspace{0.35em}
}

% --- Tighten list spacing globally ---
\setlist[itemize]{leftmargin=1.25em,itemsep=0.05em,topsep=0.1em,parsep=0pt}
\linespread{1.08}\selectfont

\setlength{\parindent}{0pt}
\pagestyle{empty}

\begin{document}

%==================== HEADER ====================
\begin{center}
  {\fontsize{28pt}{32pt}\selectfont\rmfamily\bfseries\color{accentcol} 陈冠宇}\par
  \vspace{0.9em}
  {\sffamily\fontsize{9.5pt}{12pt}\selectfont\color{subdued}\MakeUppercase{%
    Peking University \textperiodcentered\ B.E. Computer Science \textperiodcentered\ since 2025%
  }}\par
  \vspace{0.6em}
  {\rmfamily\fontsize{10pt}{12pt}\selectfont\color{subdued}{%
    \href{mailto:logic.dreemurr@gmail.com}{logic.dreemurr@gmail.com}%
    \textperiodcentered\ %
    \href{https://github.com/logic61}{github.com/logic61}%
  }}\par
\end{center}

\vspace{0.2em}
{\color{rulecol}\rule{\linewidth}{0.4pt}}
\vspace{0.4em}

%==================== I. 研究兴趣 ====================
\section{研究兴趣}
\noindent\color{bodycol}仍在持续探索.

\vspace{0.3em}

%==================== II. 教育背景 ====================
\section{教育背景}
\cventry{2022.09 -- 2025.06}{厦门外国语学校}{高中}{
  \item 第 40 届、41 届全国中学生物理竞赛（CPHO）省一等奖.
}
\cventry{2025.09 -- 至今}{北京大学}{信息科学技术学院 \textperiodcentered\ 本科 \textperiodcentered\ 计算机科学与技术}{
  \item 相关课程：程序设计实习、人工智能基础.
}

%==================== III. 项目与研究 ====================
\section{项目与研究}
\cventry{2026 --}{犬类情绪识别 \href{https://github.com/ruguo0119/dog-emotion-recognition}{\textnormal{\textit{\scriptsize (GitHub)}}}}{研究项目 \textperiodcentered\ Python \textperiodcentered\ PyTorch \textperiodcentered\ ---\ 北京大学创新+工作站}{
  \item 基于姿态与声纹特征对犬类意图进行多模态识别.
  \item 已通过开题答辩，项目持续推进中.
}
\cventry{2025}{PKU-Planner \href{https://github.com/Logic61/PKU-Planner}{\textnormal{\textit{\scriptsize (GitHub)}}}}{软件项目 \textperiodcentered\ C++ \textperiodcentered\ Qt \textperiodcentered\ --- 独立项目}{
  \item 面向北京大学教学网的课程与 DDL 桌面助手，支持教学网导入与截图识别课程.
  \item 以视觉驱动的紧迫度提示展示任务；集成一键 AI 个性化建议.
}
\cventry{2025}{Obsidian-to-Hexo \href{https://github.com/Logic61/Obsidian-to-Hexo}{\textnormal{\textit{\scriptsize (GitHub)}}}}{工具 \textperiodcentered\ Python \textperiodcentered\ --- 独立项目}{
  \item 将 Obsidian 笔记迁移至 Hexo 博客的自动化工具.
  \item 解析 Obsidian 特有 Markdown 语法，整理 frontmatter 并写入 Hexo posts 目录，支持批量备份与发布.
}

%==================== IV. 技术能力 ====================
\section{技术能力}
\cvskill{编程语言}{C++、Python}
\cvskill{框架}{PyTorch}
\cvskill{GPU 与编译器}{CUDA、Triton、TileLang（入门）}
\cvskill{工具}{Git / GitHub、LaTeX、Markdown}

%==================== V. 业余兴趣 ====================
\section{业余兴趣}
\cventry{2025.09 -- 至今}{北京大学导引队}{队员 \textperiodcentered\ 五禽戏项目代表}{}
\cventry{2026 Spring}{团体第四名 \textperiodcentered\ 五禽戏项目}{首都高校健身气功比赛 \textperiodcentered\ 代表北京大学导引队参赛}{}
\cventry{持续中}{钢琴演奏}{演奏曲目包括肖邦第一叙事曲、肖邦华丽大波兰舞曲、李斯特《弄臣》与《追雪》.}{}

%==================== FOOTER ====================
\vfill
\begin{center}
  \rmfamily\fontsize{9pt}{11pt}\selectfont\color{subdued}\MakeUppercase{%
    最后更新 2026 年 9 月 \textperiodcentered\ 中国 \textperiodcentered\ 北京%
  }
\end{center}

\end{document}
